% Extracted from the latest supplied AOS manuscript, October 29, 2026.
% Include in the manuscript with its existing macros and TikZ libraries.
\begin{figure}[ht]
\centering
\makebox[\textwidth][c]{%
    \hspace*{-0.4cm}%
\begin{tikzpicture}[
    font=\footnotesize,
    line cap=round,
    line join=round,
    >=Latex,
    manifold/.style={blue!75!black,line width=1.1pt},
    cellray/.style={black!22,line width=0.35pt},
    selectedplusray/.style={orange!85!black,line width=0.9pt},
    selectedminusray/.style={green!60!black,line width=0.9pt},
    branchdivider/.style={black!45,densely dashed,line width=0.6pt},
    guide/.style={black!48,densely dashed,line width=0.55pt},
    radiusguide/.style={purple!75!black,densely dotted,line width=0.8pt},
    plusguide/.style={
        draw=orange!95!black,densely dotted,line width=1.0pt
    },
    minusguide/.style={
        draw=green!65!black,densely dotted,line width=1.0pt
    },
    pluspoint/.style={
        circle,fill=orange!95!black,draw=white,
        line width=0.4pt,inner sep=1.8pt
    },
    minuspoint/.style={
        circle,fill=white,draw=green!65!black,
        line width=0.9pt,inner sep=1.6pt
    },
    labelbox/.style={
        fill=white,fill opacity=0.94,text opacity=1,inner sep=1.5pt
    }
]
% =========================================================
% Left panel: angular cells for q=1
% =========================================================
\begin{scope}[shift={(3.35,3.20)}]
    % Exact q=1 values from the construction:
    % eta_1   = 0.05
    % alpha_1 = 0.840971
    % beta_1  = 1.169111
    % w_1     = beta_1-alpha_1
    % kappa_1 = 10
    %
    % The radii are multiplied by 2.25 only for display.
    \def\radscale{2.25}
    \def\alphaval{0.840971}
    \def\betaval{1.169111}
    \def\midval{1.005041}
    \def\halfwidth{0.164070}
    \def\rqval{1.087076} % alpha_1 + 3w_1/4
    % Schematic display radius for the benchmark realization.
    \def\rstardisplayrad{1.68}
    \pgfmathsetmacro{\alpharad}{\radscale*\alphaval}
    \pgfmathsetmacro{\betarad}{\radscale*\betaval}
    \pgfmathsetmacro{\rqrad}{\radscale*\rqval}
    % Highlighted cell:
    % I_{1,1}=[2pi/10,4pi/10)=[36 degrees,72 degrees).
    \fill[orange!10]
        (0,0)--(36:3.02)
        arc[start angle=36,end angle=54,radius=3.02]--cycle;
    \fill[green!8]
        (0,0)--(54:3.02)
        arc[start angle=54,end angle=72,radius=3.02]--cycle;
    % Ten angular cells.
    \foreach \a in {0,36,...,324}{
        \draw[cellray] (0,0)--(\a:2.90);
    }
    % Boundaries of the selected cell.
    \draw[selectedplusray] (0,0)--(36:3.05);
    \draw[selectedminusray] (0,0)--(72:3.05);
    \draw[branchdivider] (0,0)--(54:3.05);
    % Radial guide circles.
    \draw[guide] (0,0) circle[radius=\alpharad];
    \draw[guide] (0,0) circle[radius=\betarad];
    \draw[radiusguide] (0,0) circle[radius=\rqrad];
    % Corrugated manifold:
    % rho_1(theta)
    % = alpha_1 + w_1/2 {1-cos(kappa_1 theta)}
    % = 1.005041 - 0.164070 cos(10 theta).
    \draw[manifold]
        plot[
            domain=0:360,
            samples=721,
            smooth,
            variable=\t
        ]
        ({\t}:{\radscale*(\midval-\halfwidth*cos(10*\t))});
    % Label all cells except the highlighted cell, which receives
    % a separate label below.
    \foreach \j in {0,...,9}{
        \ifnum\j=1\relax
        \else
            \node[font=\scriptsize]
                at ({36*\j+18}:3.24)
                {$\mathcal I_{\j,1}$};
        \fi
    }
    % Theta_star lies in I_{1,1}; hence J_1=1.
    % The black point is a realized noisy observation X_star.
    \draw[black,line width=0.7pt]
        (0,0)--(52:2.90);
    \fill[black]
        (52:\rstardisplayrad) circle[radius=1.8pt];
    % The selected radius has two inverse-branch locations in cell J_1=1.
    \draw[orange!90!black,line width=0.7pt]
        (0,0)--(48:\rqrad);
    % B_1=0 selects the increasing (+) branch.
    \node[pluspoint] (sone)
        at (48:\rqrad) {};
    % B_1=1 selects the decreasing (-) branch.
    \node[minuspoint] (sonealt)
        at (60:\rqrad) {};
    % The displayed coupling uses B_1=0, so the orange branch is selected.
    % Place the branch labels outside the highlighted cell and use leaders.
    \node[
        orange!85!black,
        labelbox,
        anchor=west
    ] (sonelabel)
        at (3.08,1.60) {$\boldsymbol S_1\;(B_1=0)$};
    \draw[orange!85!black,->]
        (sonelabel.west)
        to[out=180,in=-20]
        (sone);
    \node[
        green!55!black,
        labelbox,
        anchor=east
    ] (sonealtlabel)
        at (0.58,2.48) {$B_1=1$};
    \draw[green!55!black,->]
        (sonealtlabel.east)
        to[out=0,in=150]
        (sonealt);
    % Origin.
    \fill[black]
        (0,0) circle[radius=1.25pt];
    \node[anchor=north east]
        at (-0.06,-0.05)
        {$\boldsymbol 0$};
    % Selected-cell label placed outside the highlighted region.
    \node[
        anchor=south,
        orange!70!black,
        labelbox
    ]
        at (48:3.72)
        {$\mathcal I_{J_1,1},\quad J_1=1$};
    % Label the black point outside the highlighted cell.
    \node[
        anchor=west,
        align=left,
        labelbox
    ] (xstarlabel)
        at (3.08,2.30)
        {noisy benchmark realization\\
        $\boldsymbol X_\star$};
    \draw[black!75,->]
        (xstarlabel.west)
        to[out=180,in=-15]
        (52:\rstardisplayrad);
    % Radial labels placed outside the guide circles.
    \node[anchor=east,labelbox] (betalabel)
        at (-3.02,1.57)
        {$\beta_1$};
    \draw[black!55,->]
        (betalabel.east)--(154:\betarad);
    \node[
        anchor=east,
        purple!75!black,
        labelbox
    ] (rqlabel)
        at (-3.02,0.32)
        {$R_1$};
    \draw[purple!75!black,->]
        (rqlabel.east)--(171:\rqrad);
    \node[anchor=east,labelbox] (alphalabel)
        at (-3.02,-0.53)
        {$\alpha_1$};
    \draw[black!55,->]
        (alphalabel.east)--(196:\alpharad);
    \node[font=\small]
        at (0,3.82)
        {Angular cells for $q=1$ ($\kappa_1=10$)};
\end{scope}
% =========================================================
% Right panel: selected cell in local phase coordinates
% =========================================================
\begin{scope}[shift={(8.05,1.10)}]
    % Match the two phase branches to the colors in the selected cell.
    \fill[orange!8] (0,0) rectangle (3.60,2.88);
    \fill[green!7] (3.60,0) rectangle (7.20,2.88);
    % Horizontal coordinate:
    % phi=kappa_1 theta-2pi J_1 in [0,2pi].
    % The display maps [0,2pi] to [0,7.2].
    \draw[->]
        (-0.16,0)--(7.72,0);
    \draw[->]
        (0,-0.16)--(0,3.12)
        node[above] {$\rho_1$};
    \draw[branchdivider] (3.60,0)--(3.60,2.50);
    % One minimum-to-maximum-to-minimum corrugation.
    \draw[manifold]
        plot[
            domain=0:360,
            samples=241,
            smooth,
            variable=\t
        ]
        ({\t/50},{1.5-cos(\t)});
    % Purple guide for the selected radius R_1.
    \draw[radiusguide]
        (0.12,2.0)--(7.20,2.0);
    % Increasing (+) branch: phi=xi_1(R_1).
    \draw[plusguide]
        (2.40,0)--(2.40,2.0);
    % Decreasing (-) branch: phi=2pi-xi_1(R_1).
    \draw[minusguide]
        (4.80,0)--(4.80,2.0);
    \node[pluspoint]
        at (2.40,2.0) {};
    \node[minuspoint]
        at (4.80,2.0) {};
    % Branch descriptions are separated from the curve and from each other.
    \node[
        anchor=south east,
        align=right,
        labelbox,
        text=orange!95!black
    ] (pluslabel)
        at (2.12,3.00)
        {$\theta_{J_1,1}^{+}(R_1)$\\
         $B_1=0$};
    \draw[orange!95!black,->]
        (pluslabel.south east)
        to[out=-35,in=120]
        (2.40,2.05);
    \node[
        anchor=south west,
        align=left,
        labelbox,
        text=green!65!black
    ] (minuslabel)
        at (5.08,3.00)
        {$\theta_{J_1,1}^{-}(R_1)$\\
         $B_1=1$};
    \draw[green!65!black,->]
        (minuslabel.south west)
        to[out=-145,in=60]
        (4.80,2.05);
    % Radial-axis labels are placed to the left of the axis.
    \draw
        (0.08,0.5)--(-0.08,0.5);
    \node[anchor=east,labelbox]
        at (-0.13,0.5)
        {$\alpha_1$};
    \draw
        (0.08,2.0)--(-0.08,2.0);
    \node[anchor=east,labelbox,text=purple!75!black]
        at (-0.13,2.0)
        {$R_1$};
    \draw
        (0.08,2.5)--(-0.08,2.5);
    \node[anchor=east,labelbox]
        at (-0.13,2.5)
        {$\beta_1$};
    % Phase-axis ticks and labels.
    \draw
        (0,-0.07)--(0,0.07);
    \node[anchor=north]
        at (0,-0.15)
        {$0$};
    \draw[draw=orange!95!black,line width=0.9pt]
        (2.40,-0.07)--(2.40,0.07);
    \node[
        anchor=north east,
        align=right,
        text=orange!95!black
    ]
        at (2.9,-0.17)
        {$\xi_1(R_1)$};
    \draw
        (3.60,-0.07)--(3.60,0.07);
    \node[anchor=north]
        at (3.60,-0.17)
        {$\pi$};
    \draw[draw=green!65!black,line width=0.9pt]
        (4.80,-0.07)--(4.80,0.07);
    \node[
        anchor=north west,
        align=left,
        text=green!65!black
    ]
        at (4.2,-0.17)
        {$2\pi-\xi_1(R_1)$};
    \draw
        (7.20,-0.07)--(7.20,0.07);
    \node[anchor=north]
        at (7.20,-0.17)
        {$2\pi$};
    % Axis title is below all tick annotations.
    \node[anchor=north]
        at (3.60,-0.92)
        {$\phi=\kappa_1\theta-2\pi J_1$};
    % The mirrored brace opens toward the plotted radial interval.
    %\draw[
    %    decorate,
    %    decoration={brace,mirror,amplitude=4pt}
    %]
     %   (7.43,0.5)--(7.43,2.5)
      %  node[
       %     midway,
        %    right=7pt,
         %   align=left
        %]
        %{$w_1=\beta_1-\alpha_1$};
    %\node[font=\small]
    %    at (3.60,3.65)
    %    {The highlighted cell in phase coordinates};
\end{scope}
\end{tikzpicture}
    \hspace*{0.4cm}%
}
\caption{
Schematic illustration, not drawn to scale, of one realization of the
coupling between the noisy benchmark and the $q=1$ corrugated
alternative. Here $\kappa_1=10$, so
each angular cell has width $2\pi/\kappa_1=36^\circ$.
\textit{Left:} The black point is a noisy benchmark realization
$\boldsymbol X_\star=R_\star\boldsymbol u(\Theta_\star)$. Its polar angle
$\Theta_\star$ lies in the shaded cell
$\mathcal I_{1,1}=[36^\circ,72^\circ)$ and therefore determines
$J_1=1$. The orange and green halves of the cell denote its increasing
and decreasing branches, separated at $54^\circ$ (local phase $\pi$).
Compressing the radial probability coordinate of $R_\star$
gives the displayed radius
$R_1=\alpha_1+3(\beta_1-\alpha_1)/4$. The purple dotted circle of radius
$R_1$ intersects the
corrugated manifold at two eligible locations within the selected cell.
For the displayed draw $B_1=0$, the increasing-branch location
$\theta_{J_1,1}^{+}(R_1)$ is selected, giving the filled orange
latent point $\boldsymbol S_1$; the open green point is the location that
would instead be selected if $B_1=1$. The corresponding noisy
corrugated observation is subsequently obtained as
$\boldsymbol X_1=\boldsymbol S_1+\sigma_1\boldsymbol Z$ and is not shown.
\textit{Right:} The selected cell is displayed in the local phase
coordinate $\phi=\kappa_1\theta-2\pi J_1$, which maps
$\mathcal I_{J_1,1}$ onto $[0,2\pi)$. In this coordinate,
$\rho_1=\alpha_1+\{(\beta_1-\alpha_1)/2\}(1-\cos\phi)$. The orange and
green backgrounds distinguish the increasing and decreasing branches,
whose boundary is $\phi=\pi$. The colored vertical guides mark the two
phase solutions $\xi_1(R_1)$ and $2\pi-\xi_1(R_1)$, while the purple
horizontal guide marks the selected radius $R_1$.}
\label{fig:corrugation-angular-cells}
\end{figure}
